\chapter{Aksjomaty przeliczalności i własności pokrewne}
\label{ch:topog-przeliczalnosc}

\noindent\textbf{Cel rozdziału.}
Baza topologii może być ogromna nawet wtedy, gdy pojedyncze punkty
mają bardzo proste otoczenia. Rozróżnimy warunek lokalny
(pierwszy aksjomat przeliczalności) od globalnego
(drugi aksjomat przeliczalności),
a następnie porównamy je z ośrodkowością i własnością Lindelöfa.
Podamy wszystkie istotne implikacje, przeanalizujemy ograniczenia
ich odwracania oraz wykażemy, że w przestrzeniach metrycznych
trzy z tych własności są równoważne. Nie zakładamy aksjomatów
oddzielania, o ile nie zostaną wyraźnie wymienione.

\section{Pierwszy i drugi aksjomat przeliczalności}
\label{sec:topog-dwie-przeliczalnosci}

\begin{definicja}[Drugi aksjomat przeliczalności]
\label{def:topog-druga-przeliczalnosc}
Niech $(X,\tau)$ będzie przestrzenią topologiczną.
Mówimy, że $X$ \emph{spełnia drugi aksjomat przeliczalności},
jeżeli istnieje co najwyżej przeliczalna rodzina $\mathcal B\subseteq\tau$
taka, że dla każdego $U\in\tau$ i $x\in U$ istnieje
$B\in\mathcal B$ spełniający $x\in B\subseteq U$.
Innymi słowy, $\mathcal B$ jest przeliczalną bazą \emph{całej}
topologii. Dla $X=\varnothing$ można wziąć $\mathcal B=\varnothing$.
\end{definicja}

\begin{uwaga}[Warunek lokalny a globalny]
\label{rem:topog-lokalnie-globalnie-przeliczalnosc}
Pierwszy aksjomat przeliczalności zdefiniowaliśmy w rozdziale~06:
każdy punkt ma co najwyżej przeliczalną bazę otoczeń.
Rodzina tych baz może jednak zmieniać się z punktem i nie musi
istnieć jedna przeliczalna rodzina wystarczająca dla całej przestrzeni.
Bazę otoczeń rozumiemy jak w rozdziale~01: jej elementy są
otoczeniami punktu, niekoniecznie zbiorami otwartymi.
\end{uwaga}

\begin{twierdzenie}[Drugi aksjomat przeliczalności implikuje pierwszy]
\label{thm:topog-druga-pierwsza}
Każda przestrzeń spełniająca drugi aksjomat przeliczalności
spełnia pierwszy aksjomat przeliczalności.
\end{twierdzenie}
\begin{proof}
Niech $\mathcal B$ będzie co najwyżej przeliczalną bazą $X$.
Dla każdego $x\in X$ rodzina
$\mathcal B_x=\{B\in\mathcal B:x\in B\}$
jest co najwyżej przeliczalna. Jeśli $U$ jest dowolnym
otwartym otoczeniem $x$, z definicji bazy istnieje
$B\in\mathcal B_x$ z $B\subseteq U$. Jeżeli $N$ jest
otoczeniem niekoniecznie otwartym, zawiera takie $U$,
więc pewien $B\in\mathcal B_x$ zawiera $x$ i jest zawarty w $N$.
Zatem $\mathcal B_x$ jest bazą otoczeń $x$.
Dla przestrzeni pustej warunek jest spełniony automatycznie.
\end{proof}

\begin{przyklad}[Baza euklidesowa o wymiernych parametrach]
\label{ex:topog-przeliczalna-baza-euklidesowa}
Dla $n\geq1$ rodzina wszystkich kul otwartych
$B(q,r)\subseteq\R^n$ o środku $q\in\Q^n$
i promieniu $r\in\Q$, $r>0$, jest przeliczalna.
Jest bazą topologii euklidesowej: gdy $x\in U$ i $U$ jest otwarty,
istnieje $\varepsilon>0$ z $B(x,\varepsilon)\subseteq U$.
Wybieramy $q\in\Q^n$ tak blisko $x$, aby
$\|x-q\|<\varepsilon/4$, a następnie
$r\in\Q$ takie, że
$\|x-q\|<r<\varepsilon-\|x-q\|$.
Wówczas $x\in B(q,r)$, a z nierówności trójkąta
$B(q,r)\subseteq B(x,\varepsilon)\subseteq U$.
Zatem każda $\R^n$ spełnia drugi aksjomat przeliczalności.
\end{przyklad}

\begin{twierdzenie}[Dziedziczenie przeliczalności]
\label{thm:topog-przeliczalnosc-dziedziczenie}
Niech $A$ będzie podprzestrzenią przestrzeni topologicznej $X$.
Jeżeli $X$ spełnia pierwszy aksjomat przeliczalności, to $A$
również go spełnia. Jeżeli $X$ spełnia drugi aksjomat
przeliczalności, to $A$ również go spełnia.
\end{twierdzenie}
\begin{proof}
Ustalmy $a\in A$. Jeśli $\mathcal N_a$ jest przeliczalną
bazą otoczeń $a$ w $X$, to
$\{N\cap A:N\in\mathcal N_a\}$
jest bazą otoczeń $a$ w $A$. Istotnie, względne otoczenie
$a$ zawiera zbiór postaci $U\cap A$ dla pewnego otwartego
$U\subseteq X$ z $a\in U$; istnieje $N\in\mathcal N_a$
takie, że $a\in N\subseteq U$, więc
$a\in N\cap A\subseteq U\cap A$.
Jeśli $\mathcal B$ jest przeliczalną bazą $X$, to
$\mathcal B|_A=\{B\cap A:B\in\mathcal B\}$
jest przeliczalną bazą topologii podprzestrzeni:
każdy względnie otwarty zbiór ma postać $U\cap A$ i jest
sumą tych $B\cap A$, dla których $B\subseteq U$.
Oba stwierdzenia obejmują również $A=\varnothing$.
\end{proof}

\begin{twierdzenie}[Przeliczalne iloczyny przestrzeni spełniających drugi aksjomat przeliczalności]
\label{thm:topog-druga-iloczyn-przeliczalny}
Niech $I$ będzie zbiorem co najwyżej przeliczalnym,
a $\{X_i:i\in I\}$ rodziną przestrzeni spełniających
drugi aksjomat przeliczalności.
Wówczas iloczyn $\prod_{i\in I}X_i$, wyposażony w topologię
iloczynu, spełnia drugi aksjomat przeliczalności.
\end{twierdzenie}
\begin{proof}
Dla każdego $i\in I$ wybieramy co najwyżej przeliczalną bazę
$\mathcal B_i$ przestrzeni $X_i$. Dla skończonego $F\subseteq I$
i zbiorów $B_i\in\mathcal B_i$ rozważamy cylinder
\[
 C(F,(B_i))=\bigcap_{i\in F}\pi_i^{-1}(B_i).
\]
Dla $F=\varnothing$ jest to cały iloczyn. Zbiory te tworzą
bazę topologii iloczynu, ponieważ każdy podstawowy otwarty
cylinder jest takiego rodzaju po ewentualnym zastąpieniu
otwartych zbiorów w czynnikach elementami odpowiednich baz.
Rodzina skończonych podzbiorów przeliczalnego $I$ jest
przeliczalna: jest sumą po $m\geq0$ obrazów zbiorów
$I^m$ przez odwzorowanie $(i_1,\ldots,i_m)\mapsto
\{i_1,\ldots,i_m\}$. Dla ustalonego skończonego $F$
rodzina kombinacji $(B_i)_{i\in F}$ również jest
co najwyżej przeliczalna, jako skończony iloczyn rodzin
przeliczalnych. Suma przeliczalnie wielu takich rodzin jest przeliczalna.
W przypadku pustego czynnika iloczyn jest pusty,
a dla $I=\varnothing$ jest jednoelementowy; w obu sytuacjach
wniosek zachodzi bezpośrednio.
\end{proof}

\begin{przyklad}[Nieprzeliczalny iloczyn nie musi spełniać pierwszego aksjomatu przeliczalności]
\label{ex:topog-iloczyn-nieprzeliczalny}
Niech $I$ będzie nieprzeliczalny. W przestrzeni
$X=\R^I$ z topologią iloczynu oznaczmy przez $0$
punkt o wszystkich współrzędnych równych zero.
Pokażemy, że $X$ nie spełnia pierwszego aksjomatu
przeliczalności w punkcie $0$.
Załóżmy przeciwnie, że $\{U_n:n\geq1\}$ jest przeliczalną
bazą otoczeń $0$. Dla każdego $n$ istnieje podstawowy
otwarty cylinder $C_n$ z
$0\in C_n\subseteq U_n$, ograniczający tylko skończenie
wiele współrzędnych $F_n\subseteq I$. Zbiór
$F=\bigcup_{n\geq1}F_n$ jest przeliczalny, więc istnieje
$i_0\in I\setminus F$. Zbiór
$V=\pi_{i_0}^{-1}((-1,1))$ jest otwartym otoczeniem $0$.
Dla każdego $n$ cylinder $C_n$ dopuszcza dowolną wartość
współrzędnej $i_0$, ponieważ $i_0\notin F_n$.
Stąd $C_n\not\subseteq V$ i tym bardziej
$U_n\not\subseteq V$, sprzecznie z własnością bazy
otoczeń. Jest to przykład zawodności pierwszego aksjomatu
przeliczalności w dowolnych iloczynach, mimo że każdy
czynnik spełnia nawet drugi aksjomat przeliczalności.
\end{przyklad}

\section{Ośrodkowość i własność Lindelöfa}
\label{sec:topog-osrodkowosc-lindelof}

\begin{definicja}[Ośrodkowość]
\label{def:topog-osrodkowosc}
Przestrzeń topologiczna $X$ jest \emph{ośrodkowa}, jeżeli
istnieje co najwyżej przeliczalny zbiór $D\subseteq X$
gęsty w $X$, to znaczy $\overline D=X$.
Dla przestrzeni pustej można przyjąć $D=\varnothing$.
\end{definicja}

\begin{definicja}[Własność Lindelöfa]
\label{def:topog-lindelof}
Przestrzeń topologiczna $X$ jest \emph{przestrzenią Lindelöfa},
jeśli z każdego pokrycia $X$ zbiorami otwartymi można
wybrać co najwyżej przeliczalne podpokrycie.
Przestrzeń pusta spełnia ten warunek.
\end{definicja}

\begin{twierdzenie}[Dwie konsekwencje drugiego aksjomatu przeliczalności]
\label{thm:topog-druga-osrodkowa-lindelof}
Jeśli przestrzeń topologiczna $X$ spełnia drugi aksjomat
przeliczalności,
to jest ośrodkowa i jest przestrzenią Lindelöfa.
\end{twierdzenie}
\begin{proof}
Niech $\mathcal B$ będzie co najwyżej przeliczalną bazą $X$.
Dla każdego niepustego $B\in\mathcal B$ wybieramy punkt $x_B\in B$.
Zbiór $D=\{x_B:B\in\mathcal B,\ B\ne\varnothing\}$
jest co najwyżej przeliczalny. Jeśli $U\subseteq X$ jest
niepusty i otwarty, zawiera pewien niepusty
$B\in\mathcal B$, a zatem $x_B\in D\cap U$.
Każdy niepusty zbiór otwarty przecina $D$,
co jest równoważne równości $\overline D=X$.

Niech teraz $\mathcal U$ będzie pokryciem otwartym $X$.
Rozważmy rodzinę
$\mathcal B'=\{B\in\mathcal B:\exists U\in\mathcal U,
\ B\subseteq U\}$. Dla każdego $B\in\mathcal B'$
wybieramy jeden zbiór $U_B\in\mathcal U$ z $B\subseteq U_B$.
Rodzina $\{U_B:B\in\mathcal B'\}$ jest co najwyżej
przeliczalna. Pokrywa $X$: dla każdego $x\in X$
istnieje $U\in\mathcal U$ z $x\in U$, a z własności bazy
istnieje $B\in\mathcal B$ z $x\in B\subseteq U$;
wtedy $B\in\mathcal B'$ i $x\in U_B$.
Oba argumenty są poprawne również dla $X=\varnothing$.
\end{proof}

\begin{uwaga}[Aksjomat wyboru]
\label{rem:topog-wybor-przeliczalnosc}
Wybór punktu $x_B$ albo zbioru $U_B$ dla wszystkich
elementów rodziny przeliczalnej korzysta co najwyżej
z przeliczalnej postaci aksjomatu wyboru; w tym skrypcie
pracujemy w zwykłym systemie ZFC. Samo twierdzenie
Tychonoffa z rozdziału~07 wymaga w ogólności silniejszego
narzędzia. Nie utożsamiamy zatem powyższych argumentów
z dowodami zupełnie pozbawionymi wyboru.
\end{uwaga}

\begin{twierdzenie}[Zachowanie ośrodkowości i własności Lindelöfa]
\label{thm:topog-osrodk-lindelof-zachowanie}
Niech $X,Y$ będą przestrzeniami topologicznymi.
\begin{enumerate}[label=(\roman*)]
\item Jeśli $X$ jest ośrodkowa, to każda niepusta otwarta
podprzestrzeń $U\subseteq X$ jest ośrodkowa.
\item Jeśli $f:X\to Y$ jest ciągła i surjektywna, to
ośrodkowość $X$ implikuje ośrodkowość $Y$, a własność
Lindelöfa przestrzeni $X$ implikuje własność Lindelöfa
przestrzeni $Y$.
\item Jeśli $X$ jest przestrzenią Lindelöfa,
to każda domknięta podprzestrzeń $F\subseteq X$ jest
przestrzenią Lindelöfa.
\item Jeśli $X=\bigcup_{n\geq1}X_n$, gdzie każde $X_n$
(jako podprzestrzeń $X$) ma własność Lindelöfa,
to $X$ również ma własność Lindelöfa.
\end{enumerate}
\end{twierdzenie}
\begin{proof}
(i) Niech $D\subseteq X$ będzie przeliczalny i gęsty.
Dla niepustego zbioru otwartego $W$ w $U$ mamy
$W=V\cap U$ dla pewnego otwartego $V\subseteq X$.
Ponieważ $U$ jest otwarty, $W$ jest niepustym zbiorem
otwartym w $X$, a więc $W\cap D\ne\varnothing$.
Zatem $D\cap U$ jest przeliczalny i gęsty w $U$.
Przypadek $U=\varnothing$ jest równie prosty.

(ii) Jeśli $D$ jest gęsty w $X$, to dla każdego niepustego
otwartego $V\subseteq Y$ zbiór $f^{-1}(V)$ jest niepusty
(z surjektywności) i otwarty (z ciągłości). Przecina więc $D$,
a zatem $V$ przecina $f(D)$; zbiór $f(D)$ jest
przeliczalny i gęsty w $Y$. Jeśli $\mathcal V$ jest
pokryciem otwartym $Y$, to
$\{f^{-1}(V):V\in\mathcal V\}$ jest pokryciem otwartym
$X$. Wybieramy przeliczalne podpokrycie
$\{f^{-1}(V_n)\}$; wtedy $\{V_n\}$ pokrywa $Y$
z surjektywności $f$.

(iii) Niech $\mathcal W$ będzie pokryciem $F$ zbiorami
otwartymi w topologii podprzestrzeni. Każde $W\in\mathcal W$
ma postać $W=U_W\cap F$, gdzie $U_W$ jest otwarty w $X$.
Rodzina $\{U_W:W\in\mathcal W\}\cup\{X\setminus F\}$
jest otwartym pokryciem $X$, bo $F$ jest domknięty.
Istnieje jej przeliczalne podpokrycie; po pominięciu
$X\setminus F$ i przecięciu pozostałych zbiorów z $F$
otrzymujemy przeliczalne podpokrycie $F$.

(iv) Niech $\mathcal U$ będzie otwartym pokryciem $X$.
Dla każdego $n$ rodzina
$\{U\cap X_n:U\in\mathcal U\}$ pokrywa $X_n$.
Wybieramy przeliczalną podrodzinę $\mathcal U_n\subseteq\mathcal U$
pokrywającą $X_n$. Suma mnogościowa
$\bigcup_{n\geq1}\mathcal U_n$ jest przeliczalna
i pokrywa $\bigcup_nX_n=X$.
\end{proof}

\begin{przyklad}[Przestrzeń ośrodkowa nie musi spełniać pierwszego aksjomatu przeliczalności]
\label{ex:topog-kosk-onczona-nie-pierwsza}
Niech $X$ będzie nieprzeliczalnym zbiorem z topologią
koskończoną. Dowolny nieskończony przeliczalny
$D\subseteq X$ jest gęsty: każdy niepusty zbiór otwarty
ma skończone dopełnienie, więc przecina $D$.
Przestrzeń jest także zwarta: z danego pokrycia otwartego
wybieramy jeden niepusty zbiór otwarty $U$, a pozostały
skończony zbiór $X\setminus U$ pokrywamy skończenie
wieloma elementami pokrycia.
Jednak $X$ nie spełnia pierwszego aksjomatu przeliczalności.
Istotnie, jeśli $\{N_n:n\geq1\}$ byłaby bazą otoczeń
punktu $x$, wybierzmy dla każdego $n$ otwarty zbiór
$U_n$ z $x\in U_n\subseteq N_n$. Zbiór
$F=\bigcup_n(X\setminus U_n)$ jest przeliczalny.
Wybierzmy $y\in X\setminus(F\cup\{x\})$. Wtedy
otwarte otoczenie $X\setminus\{y\}$ punktu $x$
nie zawiera żadnego $N_n$, gdyż
$y\in U_n\subseteq N_n$ dla każdego $n$.
Otrzymujemy sprzeczność z definicją bazy otoczeń.
W szczególności ośrodkowość, a nawet zwartość i aksjomat
$T_1$, nie zapewniają spełnienia pierwszego aksjomatu przeliczalności.
\end{przyklad}

\section{Przestrzenie metryczne: trzy równoważne własności}
\label{sec:topog-przeliczalnosc-metryczna}

\begin{twierdzenie}[Ośrodkowość, drugi aksjomat przeliczalności i własność Lindelöfa]
\label{thm:topog-metryczne-trzy-warunki}
Niech $(X,d)$ będzie przestrzenią metryczną.
Następujące warunki są równoważne:
\begin{enumerate}[label=(\roman*)]
\item $X$ jest ośrodkowa;
\item $X$ spełnia drugi aksjomat przeliczalności;
\item $X$ jest przestrzenią Lindelöfa.
\end{enumerate}
\end{twierdzenie}
\begin{proof}
Implikacje (ii)$\Rightarrow$(i) i (ii)$\Rightarrow$(iii)
wynikają z twierdzenia~\ref{thm:topog-druga-osrodkowa-lindelof}.

(i)$\Rightarrow$(ii). Niech $D\subseteq X$ będzie
co najwyżej przeliczalny i gęsty. Rodzina
\[
 \mathcal B=\{B_d(a,r):a\in D,\ r\in\Q,\ r>0\}
\]
jest co najwyżej przeliczalna. Wykażemy, że jest bazą.
Niech $x\in U$, gdzie $U$ jest otwarty. Wybierzmy
$\varepsilon>0$ z $B_d(x,\varepsilon)\subseteq U$.
Z gęstości $D$ istnieje $a\in D$ z
$d(x,a)<\varepsilon/4$. Wybieramy liczbę wymierną
$r$ spełniającą
$d(x,a)<r<\varepsilon-d(x,a)$; jest to możliwe,
ponieważ $2d(x,a)<\varepsilon/2<\varepsilon$.
Wtedy $x\in B_d(a,r)$, a jeżeli $y\in B_d(a,r)$,
to z nierówności trójkąta
$d(x,y)\leq d(x,a)+d(a,y)<d(x,a)+r<\varepsilon$.
Zatem $B_d(a,r)\subseteq U$.

(iii)$\Rightarrow$(i). Dla każdego $n\geq1$ rodzina
$\{B_d(x,1/n):x\in X\}$ jest otwartym pokryciem $X$.
Z własności Lindelöfa wybieramy co najwyżej przeliczalne
podpokrycie o środkach ze zbioru $D_n\subseteq X$.
Niech $D=\bigcup_{n\geq1}D_n$. Zbiór $D$ jest
co najwyżej przeliczalny. Dla dowolnego niepustego
zbioru otwartego $U$ wybieramy $x\in U$ i
$\varepsilon>0$ z $B_d(x,\varepsilon)\subseteq U$.
Bierzemy $n$ tak duże, że $1/n<\varepsilon$.
Ponieważ kule o środkach z $D_n$ pokrywają $X$,
istnieje $a\in D_n$ z $d(x,a)<1/n<\varepsilon$.
Wtedy $a\in D\cap U$. Zatem $D$ jest gęsty.
Jeśli $X=\varnothing$, wszystkie trzy własności
zachodzą bezpośrednio; argument można zastosować do
$X\ne\varnothing$.
\end{proof}

\begin{wniosek}[Przykłady metryczne]
\label{cor:topog-osrodkowa-metryczna-baza}
Dla każdego $n\geq1$ przestrzeń $\R^n$ z topologią
euklidesową spełnia drugi aksjomat przeliczalności,
jest ośrodkowa i ma
własność Lindelöfa. Każda jej podprzestrzeń ma te
trzy własności.
\end{wniosek}
\begin{proof}
Rodzina $\Q^n$ jest przeliczalna i gęsta w $\R^n$,
więc stosujemy twierdzenie
\ref{thm:topog-metryczne-trzy-warunki}. Dowolna podprzestrzeń
$\R^n$ jest metryczna, a drugi aksjomat przeliczalności
przechodzi na podprzestrzenie na mocy twierdzenia
\ref{thm:topog-przeliczalnosc-dziedziczenie}.
Ponownie stosując równoważność, otrzymujemy ośrodkowość
i własność Lindelöfa podprzestrzeni.
\end{proof}

\section{Prosta Sorgenfreya: ostrzeżenie przed odwracaniem implikacji}
\label{sec:topog-sorgenfrey-przeliczalnosc}

\begin{definicja}[Prosta Sorgenfreya]
\label{def:topog-sorgenfrey-linia}
Na $\R$ rozważamy topologię mającą za bazę wszystkie
półotwarte przedziały $[a,b)$, gdzie $a<b$.
Otrzymaną przestrzeń oznaczamy przez $S$ i nazywamy
\emph{prostą Sorgenfreya}. Baza ta rzeczywiście wyznacza
topologię: każdy punkt należy do pewnego takiego przedziału,
a jeśli $x\in[a,b)\cap[c,d)$, to
$[x,\min\{b,d\})$ jest zawarty w przecięciu.
\end{definicja}

\begin{figure}[htbp]
\centering
\begin{tikzpicture}[x=1cm,y=1cm]
\draw[-{Latex[length=2mm]},gray!75] (-3.4,.9)--(3.4,.9);
\draw[line width=2.1pt,blue!65!black] (-1.6,.9)--(1.65,.9);
\fill[blue!65!black] (-1.6,.9) circle (2.9pt);
\draw[blue!65!black,fill=white,line width=.95pt] (1.65,.9) circle (2.9pt);
\node[above=3mm] at (0,.9) {$[a,b)$};
\node[below=2mm] at (-1.6,.9) {$a$};
\node[below=2mm] at (1.65,.9) {$b$};
\draw[-{Latex[length=2mm]},gray!75] (-3.4,-.45)--(3.4,-.45);
\draw[line width=2.1pt,orange!80!black] (-1.6,-.45)--(1.65,-.45);
\draw[orange!80!black,fill=white,line width=.95pt] (-1.6,-.45) circle (2.9pt);
\draw[orange!80!black,fill=white,line width=.95pt] (1.65,-.45) circle (2.9pt);
\node[above=3mm] at (0,-.45) {$(a,b)$};
\node[below=2mm] at (-1.6,-.45) {$a$};
\node[below=2mm] at (1.65,-.45) {$b$};
\end{tikzpicture}
\caption{Przedział bazowy topologii Sorgenfreya ma domknięty lewy
koniec i otwarty prawy koniec (u góry). W topologii euklidesowej
przedział bazowy jest otwarty z obu stron (u dołu).}
\label{fig:topog-sorgenfrey-porownanie}
\end{figure}

\begin{twierdzenie}[Cztery własności prostej Sorgenfreya]
\label{thm:topog-sorgenfrey-wlasnosci-przeliczalnosc}
Prosta Sorgenfreya $S$ spełnia pierwszy aksjomat
przeliczalności, jest ośrodkowa i Lindelöfa,
ale nie spełnia drugiego aksjomatu przeliczalności.
\end{twierdzenie}
\begin{proof}
\emph{Pierwszy aksjomat przeliczalności.} Dla każdego $x\in\R$ rodzina
$\{[x,x+1/n):n\geq1\}$ jest przeliczalną bazą otoczeń $x$.
Rzeczywiście, każde otwarte otoczenie $x$ zawiera przedział
bazowy $[a,b)$ z $a\leq x<b$, a wówczas dla dostatecznie
dużego $n$ zachodzi $[x,x+1/n)\subseteq[a,b)$.

\emph{Ośrodkowość.} Zbiór $\Q$ jest przeliczalny i gęsty,
ponieważ każdy niepusty przedział $[a,b)$ zawiera
liczbę wymierną z $(a,b)$.

\emph{Niespełnianie drugiego aksjomatu przeliczalności.} Gdyby
$\mathcal B=\{B_n:n\geq1\}$ była przeliczalną bazą,
dla każdego $x\in\R$ istniałby najmniejszy indeks
\[
 n(x)=\min\{n:x\in B_n\subseteq[x,x+1)\}.
\]
Indeksy $n(x)$ dla różnych $x$ byłyby różne. Istotnie,
przy $x<y$ równość $B_{n(x)}=B_{n(y)}$ byłaby
niemożliwa, bo ten sam zbiór musiałby zawierać $x$,
a jednocześnie być zawarty w $[y,y+1)$,
który nie zawiera $x$. W ten sposób otrzymalibyśmy
iniekcję $\R\to\mathbb N$, sprzeczną z
nieprzeliczalnością $\R$.

\emph{Własność Lindelöfa.} Niech $\mathcal U$ będzie
dowolnym otwartym pokryciem $S$. Rozważmy rodzinę
\[
 \mathcal I=\{[a,b):a<b,\ [a,b)\subseteq U
 \text{ dla pewnego }U\in\mathcal U\}.
\]
Jest to pokrycie $\R$ przedziałami bazowymi.
Zbiór $G=\bigcup_{[a,b)\in\mathcal I}(a,b)$ jest
otwarty w zwykłej topologii euklidesowej. Ponieważ
$\R$ z tą topologią spełnia drugi aksjomat
przeliczalności, a jej podprzestrzeń $G$ również go spełnia
(twierdzenie~\ref{thm:topog-przeliczalnosc-dziedziczenie}),
a więc jest przestrzenią Lindelöfa. Z pokrycia $G$ przedziałami $(a,b)$,
gdzie $[a,b)\in\mathcal I$, możemy wybrać
przeliczalną podrodzinę $(a_n,b_n)$ pokrywającą $G$.
Dla każdego $n$ wybieramy $U_n\in\mathcal U$
z $[a_n,b_n)\subseteq U_n$. Rodzina $\{U_n\}$ pokrywa $G$.

Pozostaje zbiór $F=\R\setminus G$. Dla każdego $x\in F$
istnieje $[x,b)\in\mathcal I$ dla pewnego $b>x$.
Istotnie, dowolny przedział $[a,b)\in\mathcal I$
zawierający $x$ musi mieć $a=x$, gdyż w przypadku
$a<x<b$ mielibyśmy $x\in(a,b)\subseteq G$.
Ustalmy raz na zawsze numerację $\Q=\{q_1,q_2,\ldots\}$.
Dla każdego $x\in F$ istnieje wymierne $q>x$ takie,
że $[x,q)\in\mathcal I$: wystarczy wziąć
$q\in\Q\cap(x,b)$ i zauważyć, że $[x,q)\subseteq[x,b)$.
Niech $q_x$ będzie pierwszą taką liczbą wymierną
w ustalonej numeracji. Wówczas $(x,q_x)\subseteq G$.
Jeśli $x,y\in F$ oraz $x<y$, to $y\notin(x,q_x)$,
a zatem $q_x\leq y<q_y$. Stąd odwzorowanie
$x\mapsto q_x$ jest iniekcją $F\to\Q$, więc $F$
jest co najwyżej przeliczalny. Dla każdego $x\in F$
wybieramy $U_x\in\mathcal U$ zawierający $[x,q_x)$,
co jest możliwe z definicji $\mathcal I$.
Przeliczalna rodzina $\{U_n\}\cup\{U_x:x\in F\}$
pokrywa zarówno $G$, jak i $F$, a zatem całe $S$.
Dowodzi to własności Lindelöfa.
\end{proof}

\begin{wniosek}[Trzy warunki nie zapewniają drugiego aksjomatu przeliczalności]
\label{cor:topog-sorgenfrey-trzy-nie-druga}
Pierwszy aksjomat przeliczalności, ośrodkowość
i własność Lindelöfa, nawet przy założeniu Hausdorffa,
nie implikują łącznie drugiego aksjomatu przeliczalności.
\end{wniosek}
\begin{proof}
Prosta Sorgenfreya spełnia trzy wymienione warunki,
lecz nie spełnia drugiego aksjomatu przeliczalności.
Jest Hausdorffa:
jeśli $x<y$, to przedziały $[x,y)$ i $[y,y+1)$
są rozłącznymi otwartymi otoczeniami odpowiednio $x$ i $y$.
\end{proof}

\begin{przyklad}[Iloczyn przestrzeni Lindelöfa nie musi być przestrzenią Lindelöfa]
\label{ex:topog-sorgenfrey-iloczyn-nie-lindelof}
Płaszczyzna Sorgenfreya $S\times S$ jest iloczynem
przestrzeni Lindelöfa, lecz sama nie ma własności Lindelöfa.
Uzasadnimy to na podstawie zbioru
$D=\{(t,-t):t\in\R\}$, który pojawił się również
w dowodzie twierdzenia o płaszczyźnie Sorgenfreya
w rozdziale~05.
Dla $(x,-x)\in D$ podstawowy prostokąt
$[x,x+\varepsilon)\times[-x,-x+\varepsilon)$
przecina $D$ tylko w $(x,-x)$: jeśli $(t,-t)$ należy
do prostokąta, mamy jednocześnie $t\ge x$ i $t\le x$.
Zatem $D$ jest podprzestrzenią dyskretną.
Jest też domknięty. Jeśli $(u,v)\notin D$, to $u+v\ne0$.
Dla $u+v>0$ każdy podstawowy prostokąt wychodzący
z $(u,v)$ składa się z punktów o dodatniej sumie
współrzędnych; dla $u+v<0$ wystarczy wybrać
$\varepsilon>0$ tak małe, aby $u+v+2\varepsilon<0$.
W obu przypadkach otrzymujemy otwarte otoczenie
nieprzecinające $D$.

Nieprzeliczalna przestrzeń dyskretna nie jest przestrzenią
Lindelöfa, ponieważ pokrycie wszystkimi singletonami
nie ma przeliczalnego podpokrycia. Gdyby $S\times S$
była przestrzenią Lindelöfa, jej domknięta podprzestrzeń
$D$ byłaby Lindelöfa na mocy twierdzenia
\ref{thm:topog-osrodk-lindelof-zachowanie}(iii), sprzeczność.
\end{przyklad}

\section{Które własności przechodzą na podprzestrzenie?}
\label{sec:topog-przeliczalnosc-kontrprzyklady}

\begin{przyklad}[Ośrodkowość nie przechodzi na domknięte podprzestrzenie]
\label{ex:topog-osrodkowosc-nie-dziedziczna}
Płaszczyzna Sorgenfreya $S\times S$ jest ośrodkowa:
$\Q\times\Q$ jest przeliczalny i przecina każdy
niepusty podstawowy prostokąt
$[a,b)\times[c,d)$, ponieważ każdy z obu przedziałów
zawiera liczbę wymierną. Jej domknięta podprzestrzeń
$D=\{(t,-t):t\in\R\}$ jest jednak nieprzeliczalna
i dyskretna, jak wykazaliśmy w przykładzie
\ref{ex:topog-sorgenfrey-iloczyn-nie-lindelof}.
W każdej przestrzeni dyskretnej jedynym gęstym podzbiorem
jest cała przestrzeń, ponieważ każdy singleton jest otwarty.
Zatem $D$ nie jest ośrodkowa.
\end{przyklad}

\begin{przyklad}[Przestrzeń zwarta z otwartą podprzestrzenią
niebędącą przestrzenią Lindelöfa]
\label{ex:topog-lindelof-nie-dziedziczna}
Niech $E$ będzie zbiorem nieprzeliczalnym i niech
$X=E\cup\{\infty\}$, gdzie $\infty\notin E$.
Przyjmijmy za otwarte wszystkie podzbiory $E$ oraz
zbiory zawierające $\infty$, których dopełnienie w $E$
jest skończone. Rodzina ta jest topologią:
sumy zbiorów otwartych są otwarte, a przecięcie
skończenie wielu otoczeń $\infty$ nadal ma skończone
dopełnienie w $E$.

Przestrzeń $X$ jest zwarta. Każde pokrycie otwarte ma
element $U$ zawierający $\infty$; pozostała część
$X\setminus U$ jest skończona, więc wystarczy dobrać
do $U$ skończenie wiele dalszych elementów pokrycia.
Jest też Hausdorffa: różne punkty $e,f\in E$
oddzielają singletony, a $e\in E$ i $\infty$
oddzielają zbiory $\{e\}$ oraz $X\setminus\{e\}$.

Podprzestrzeń $E$ jest otwarta i dyskretna, a ponieważ
jest nieprzeliczalna, nie jest przestrzenią Lindelöfa.
Zatem własność Lindelöfa nie przechodzi na dowolne
otwarte podprzestrzenie, mimo że przechodzi na domknięte.
\end{przyklad}

\begin{samepage}
\begin{przyklad}[Pierwszy aksjomat przeliczalności nie zapewnia ośrodkowości]
\label{ex:topog-dyskretna-nie-osrodkowa}
Każda przestrzeń dyskretna spełnia pierwszy aksjomat
przeliczalności:
dla punktu $x$ jednoelementowa rodzina $\{\{x\}\}$
jest bazą jego otoczeń. Jeśli jednak zbiór punktów $E$
jest nieprzeliczalny, topologia dyskretna na $E$ nie jest
ośrodkowa (gęsty jest jedynie cały $E$), nie jest
przestrzenią Lindelöfa (pokrycie singletonami) ani
nie spełnia drugiego aksjomatu przeliczalności
(każda baza musi zawierać
wszystkie singletony).
\end{przyklad}
\end{samepage}

\section{Zestawienie implikacji i ich ograniczeń}
\label{sec:topog-przeliczalnosc-podsumowanie}

\begin{figure}[htbp]
\centering
\begin{tikzpicture}[
 >=Stealth,
 every node/.style={font=\small,align=center},
 main/.style={draw=blue!65!black,fill=blue!6,rounded corners=2pt,
 text width=3.25cm,minimum height=9mm},
 aux/.style={draw=black!55,fill=black!3,rounded corners=2pt,
 text width=3.25cm,minimum height=9mm}]
\node[main] (second) at (0,2.1) {drugi aksjomat\\przeliczalności};
\node[aux] (first) at (-4.1,.1) {pierwszy aksjomat\\przeliczalności};
\node[aux] (separ) at (0,.1) {ośrodkowość};
\node[aux] (lind) at (4.1,.1) {własność Lindelöfa};
\draw[->,thick] (second.south west) -- (first.north east);
\draw[->,thick] (second.south) -- (separ.north);
\draw[->,thick] (second.south east) -- (lind.north west);
\end{tikzpicture}
\caption{Implikacje prawdziwe dla wszystkich przestrzeni topologicznych.
Strzałki w ogólności nie są odwracalne, nawet przy dodatkowych
warunkach podanych w tekście. Równoważność metryczna dotyczy
wyłącznie drugiego aksjomatu przeliczalności, ośrodkowości
i własności Lindelöfa.}
\label{fig:topog-przeliczalnosc-implikacje}
\end{figure}

\begin{stwierdzenie}[Implikacje bez aksjomatów oddzielania]
\label{prop:topog-przeliczalnosc-hierarchia}
Dla dowolnej przestrzeni topologicznej zachodzą implikacje
\[
\text{drugi aksjomat przeliczalności}\ \Longrightarrow\
\begin{cases}
\text{pierwszy aksjomat przeliczalności},\\
\text{ośrodkowość},\\
\text{własność Lindelöfa}.
\end{cases}
\]
Żadna z implikacji odwrotnych nie zachodzi bez dodatkowych
założeń. Ponadto żadne dwa spośród trzech warunków po prawej
nie są, w ogólności, równoważne.
\end{stwierdzenie}
\begin{proof}
Implikacje udowodniono w twierdzeniach
\ref{thm:topog-druga-pierwsza} oraz
\ref{thm:topog-druga-osrodkowa-lindelof}.
Nieprzeliczalna przestrzeń dyskretna spełnia pierwszy
aksjomat przeliczalności, lecz nie jest ośrodkowa ani Lindelöfa.
Nieprzeliczalna przestrzeń koskończona jest ośrodkowa
i Lindelöfa, ale nie spełnia pierwszego aksjomatu przeliczalności.
Prosta Sorgenfreya spełnia pierwszy aksjomat
przeliczalności, jest ośrodkowa i Lindelöfa, lecz
nie spełnia drugiego aksjomatu przeliczalności.
Natomiast płaszczyzna Sorgenfreya jest ośrodkowa
i spełnia pierwszy aksjomat przeliczalności, lecz nie jest Lindelöfa.
Z drugiej strony, zwarta przestrzeń $E\cup\{\infty\}$
z przykładu~\ref{ex:topog-lindelof-nie-dziedziczna}
nie jest ośrodkowa, ponieważ jej nieprzeliczalna otwarta
podprzestrzeń dyskretna $E$ nie może zawierać gęstego
przeliczalnego zbioru: dowolny gęsty podzbiór całego
$E\cup\{\infty\}$ musi zawierać wszystkie punkty izolowane $E$.
Przykłady te weryfikują w szczególności, że wszystkie
odwrócenia trzech głównych implikacji zawodzą,
a warunki po prawej nie są parami równoważne.
\end{proof}

\begin{uwaga}[Co następuje dalej?]
\label{rem:topog-przeliczalnosc-metryzowalnosc}
W przestrzeniach metrycznych pierwszy aksjomat przeliczalności
jest zawsze spełniony, a trzy pozostałe własności są równoważne.
W ogólnych przestrzeniach rozdział ten pokazuje już,
że nawet łączne spełnienie pierwszego aksjomatu przeliczalności,
ośrodkowości i własności Lindelöfa nie wystarcza do
metryzowalności. W szczególności prosta Sorgenfreya
nie jest metryzowalna: gdyby jej topologię wyznaczała
metryka, z ośrodkowości i twierdzenia
\ref{thm:topog-metryczne-trzy-warunki}
wynikałoby spełnienie drugiego aksjomatu przeliczalności,
co jest niemożliwe.
Pełniejsze kryteria metryzowalności należą do dalszej
części skryptu.
\end{uwaga}
\FloatBarrier
